
\section{Discussion}
\label{sec:discussion}

\paragraph{Summary of findings.}
The simulation study establishes that the Single Quantile Sampler is a competitive Bayesian estimator for quantile trend filtering across a range of signal structures, noise distributions, and quantile levels. It consistently improves on \citet{onizuka2024fast} under the Laplace prior in both MSE and MAD, and matches or approaches the horseshoe variant at interior quantile levels under all three noise distributions. The advantage of the globally penalized P-MALA approach is most pronounced on the smooth signal, where the global prior's implicit bias toward a homogeneous degree of smoothness is well matched to the truth and the sampler attains the lowest errors at nearly every interior level. The ordering reverses, as expected, on the piecewise constant signal under Gaussian and mixed normal noise at interior quantiles, where the per-coefficient horseshoe prior's local adaptivity gives it the edge. The margins in that setting are small in absolute terms---all three methods recover the signal accurately---but their direction is systematic, consistent with the general principle that local adaptivity benefits estimation when the target signal has genuinely heterogeneous smoothness, with the associated computational and modeling complexity being the price paid for that benefit.

\paragraph{Scalability.}
A central motivation for the proposed approach is to make Bayesian inference practical while retaining the optimization structure of frequentist quantile trend filtering. The supplied Gibbs implementation of \citet{onizuka2024fast} forms a dense inverse and uses dense Gaussian simulation at each iteration. Our audit found that its sparse option still enters this dense computational path in the benchmark environment, so enabling that option does not remove the principal bottleneck. The augmented hierarchy also updates observation-level latent variables and local prior scales, whereas our formulation samples one global regularization parameter and one likelihood scale per quantile in addition to the trend.

This distinction should be separated from an asymptotic claim about all Gibbs implementations. At fixed difference order, a Gaussian draw based directly on the banded precision can be computed in $O(n)$ work without forming an inverse, although the changing weights generally require a new numerical factorization each iteration. For our constrained proximal system, $\mathbf{A}=\gamma^{-1}\mathbf{I}+\rho D^\top D+\rho\mathbf{I}$ is independent of the sampled parameters, allowing factor reuse when $D$, $\gamma$, and $\rho$ remain fixed. Sparse operations give $O(nJ)$ work for $J$ inner ADMM sweeps, plus linear setup and sufficient-statistic calculations. Practical comparisons must account for the inner iteration count, factorization reuse in the actual implementation, and mixing of the complete chain. Timing per outer iteration alone does not establish cost per effective posterior draw.
% Writing note: place detailed timings and the sparse-option dispatch audit in the dedicated computational section. The existing timing report used an older PMALA proposal with a documented proposal-density mismatch; rerun validated current kernels before interpreting these timings as posterior-efficiency evidence.

\paragraph{Role of the scale parameter.}
Treating the ALD scale $\sigma$ as an unknown parameter rather than fixing it at the optimization convention of unity is a substantive modeling decision with several consequences. Within the Gibbs scheme, $\sigma$ and $\lambda$ interact with the trend estimate in complementary ways: $\sigma$ calibrates the effective likelihood temperature, controlling how sharply the gradient of the smoothed check loss responds to residuals, while $\lambda$ calibrates the effective smoothness penalty. Decoupling these two sources of regularization allows each to adapt to a different aspect of the data --- noise level and signal roughness respectively --- and prevents the single tuning parameter of the frequentist formulation from having to absorb both roles simultaneously. Learning $\sigma$ allows the likelihood scale to respond to the data, but does not by itself guarantee frequentist coverage when the asymmetric Laplace likelihood is misspecified. Posterior sampling avoids the additional mean-field approximation used by VB; coverage remains a separate question, as recognized by the calibration approach of \citet{onizuka2024fast}. A practical caveat is that $\sigma$ and $\lambda$ are not fully separable in their effects on the posterior for $\theta$: both appear in the effective curvature seen by the P-MALA proposal, and under severe misspecification their joint posterior may exhibit a soft ridge. The weakly informative prior $\sigma \sim \mathrm{IG}(0.1, 0.1)$ is sufficient to regularize this in all settings considered, but sensitivity of the posterior to this prior choice warrants attention in applications where the noise distribution is far from symmetric or unimodal.

\paragraph{Limitations.}
The globally penalized prior is the principal structural limitation of the proposed approach. By placing a single shared penalty $\lambda$ on all entries of $D\theta$, the model encodes a homogeneous prior belief about the smoothness of the trend. This is appropriate when the signal is globally smooth or when the difference order $k$ is chosen to match the signal class, as in the varying smoothness experiments, but it is a deliberate simplification relative to the per-coefficient shrinkage of the horseshoe prior. On piecewise constant signals, the globally penalized model must use a single $\lambda$ to accommodate both flat segments and jump locations, a bias-variance tradeoff that the horseshoe prior resolves more cleanly through local adaptivity. The edge of Onizuka (Horseshoe) on the PC signal under Gaussian and mixed normal noise is a direct consequence of this structural difference: it is systematic and should not be expected to vanish with larger samples or better tuning, though in the settings considered here it amounts to small absolute differences between estimators that both recover the signal well. Users whose signals are known to be piecewise constant, and for whom these margins matter, may prefer the locally adaptive prior at its higher computational cost.

\paragraph{The sentinel strategy as a heuristic.}
The sentinel quantile construction introduced in Section~\ref{subsec:sentinel} resolves the gradient degeneracy of P-MALA at extreme quantile levels and produces stable estimates in the simulation study. It is, however, a practical device rather than a formally justified procedure. The specific recommendations --- up to five sentinels per side, equally spaced within each side, with spacings chosen to fill the available probability room --- are motivated by empirical experience and by the requirement that the sentinels be close enough to the target to provide a tight dynamic bound while remaining far enough apart to avoid degeneracy among themselves. 

\paragraph{Extensions.}
Several directions for extension follow naturally from the present work. The P-MALA framework accommodates non-Euclidean preconditioning through the proximal operator, and a preconditioned variant exploiting the banded structure of $D^\top D$ may improve mixing at higher difference orders without increasing per-iteration cost. The sentinel strategy is applicable to any gradient-based sampler for quantile estimation and could be combined with locally adaptive priors, potentially recovering the accuracy advantages of the horseshoe prior at extreme quantile levels while retaining the computational structure of the present approach. Extension to multivariate or spatially indexed signals requires generalization of the difference operator $D$ to non-chain graph structures, for which the banded linear system no longer applies directly; spectral or multigrid solvers may serve a similar role in those settings. Finally, the theoretical properties of the P-MALA approximation under the non-smooth ALD likelihood --- including mixing time bounds and the effect of the smoothing parameter $\varepsilon$ on the approximation error in the acceptance ratio --- remain to be established.
